%% =============================================================================
%% Section 2 — Research Methodology
%% "How Do LLM-Based Agents Represent Code Repositories?"
%% ACM Transactions on Software Engineering and Methodology (TOSEM)
%% =============================================================================

\section{Research Methodology}
\label{sec:methodology}

Our review protocol is grounded in five complementary methodological pillars:
(1) Kitchenham and Charters' guidelines for systematic literature reviews in
software engineering~\cite{kitchenham2007guidelines};
(2) Wohlin's snowballing procedure for extending database
searches~\cite{wohlin2014snowballing};
(3) the PRISMA 2020 reporting standard for systematic
reviews~\cite{page2021prisma};
(4) Cruzes and Dyb{\aa}'s recommended steps for thematic synthesis in
software engineering~\cite{cruzes2011recommended};
and (5) the ACM SIGSOFT Empirical Standards for secondary
studies~\cite{acmsigsoft2020standards}.
The integration of these five pillars is unusual — none of the five surveys
reviewed in Section~\ref{sec:relatedwork} applies more than two — and it
directly affects the completeness, reproducibility, and analytical depth of
the review.

\subsection{Research Questions}
\label{subsec:rqs}

\begin{description}
  \item[\textbf{RQ1.}] What paradigms do LLM-based systems employ to represent
        a code repository as context?
        (\emph{Taxonomy construction})

  \item[\textbf{RQ2.}] Along which structural dimensions do the seven paradigms
        differ?
        (\emph{Comparative analysis})

  \item[\textbf{RQ3.}] How are repository representation paradigms distributed
        across seven SE objectives?
        (\emph{Cross-task mapping})

  \item[\textbf{RQ4.}] What open challenges and research directions does
        the synthesis reveal?
        (\emph{Research agenda})
\end{description}

\subsection{Search Strategy}
\label{subsec:search}

We searched four primary bibliographic databases (IEEE Xplore, ACM Digital
Library, Springer Link, arXiv) using the following query, applied to title,
abstract, and keywords:

\begin{quote}
\texttt{("repository-level" OR "repository level" OR "cross-file" OR
"repo-level") AND ("large language model" OR "LLM" OR "code LLM") AND
("code completion" OR "code generation" OR "program repair" OR "bug
localization" OR "code search" OR "test generation" OR "code translation")}
\end{quote}

The search was conducted in January 2026, covering publications from
January 2020 to January 2026.
The temporal lower bound was chosen to capture the first LLM-era work on
repository-level SE (CodeBERT, 2020~\cite{feng2020codebert}).

\paragraph{Snowballing.}
Following Wohlin~\cite{wohlin2014snowballing}, we performed one round of
backward snowballing (scanning references of included papers) and one round of
forward snowballing (scanning papers that cite included papers via Semantic
Scholar and Google Scholar).
The snowballing round extended the candidate pool by 31 additional papers,
of which 19 passed the eligibility criteria.

\subsection{Inclusion and Exclusion Criteria}
\label{subsec:criteria}

\paragraph{Inclusion criteria.}

\begin{description}
  \item[\textbf{IC1.}] The paper addresses a repository-level or cross-file
        SE task, i.e., the input to the LLM contains information from more
        than one source file or requires reasoning about inter-file
        dependencies.

  \item[\textbf{IC2.}] The paper explicitly describes, proposes, or evaluates
        a strategy for constructing or selecting LLM context from a code
        repository.

  \item[\textbf{IC3.}] The paper is published in or after 2020 and is written
        in English.

  \item[\textbf{IC4.}] The paper is a full research paper (workshop papers,
        position papers, and tool demonstrations are included if they satisfy
        IC1 and IC2 and contain empirical evaluation).
\end{description}

\paragraph{Exclusion criteria.}

\begin{description}
  \item[\textbf{EC1.}] Papers whose \emph{primary} contribution is a new
        fine-tuning dataset or a fine-tuning technique, where context
        construction is not the central concern.

  \item[\textbf{EC2.}] Papers that study function-level or file-level code
        tasks without any cross-file component.

  \item[\textbf{EC3.}] Duplicate publications (conference version vs.\ extended
        journal version: the journal version is retained).

  \item[\textbf{EC4.}] Secondary studies, surveys, and literature reviews
        (these are discussed separately in Section~\ref{sec:relatedwork}).
\end{description}

\subsection{Study Selection Process and PRISMA Flow}
\label{subsec:prisma}

Figure~\ref{fig:prisma} shows the PRISMA 2020 flow
diagram~\cite{page2021prisma} documenting every stage of the selection process.
The complete PRISMA 2020 checklist is provided in Online Supplement~A.

\begin{figure}[t]
  \centering
  % PRISMA 2020 flow diagram
  \begin{tikzpicture}[
      box/.style={rectangle, draw, rounded corners=4pt,
                  text width=5cm, minimum height=1.0cm,
                  align=center, font=\small},
      eliminated/.style={rectangle, draw, dashed, rounded corners=4pt,
                  text width=4.2cm, minimum height=0.8cm,
                  align=center, font=\scriptsize, text=gray},
      arr/.style={->, >=Stealth, thick},
    ]
    % Identification
    \node[box] (db)    at (0,0)    {Database search ($n = 1{,}847$)};
    \node[box] (snow)  at (0,-1.8) {Snowballing ($n = 31$)};
    \node[box] (dedup) at (0,-3.6) {After deduplication ($n = 1{,}623$)};
    % Screening
    \node[box]       (screen) at (0,-5.4)   {Title \& abstract screening ($n = 1{,}623$)};
    \node[eliminated](excl1)  at (5.8,-5.4) {Excluded on title/abstract ($n = 1{,}127$)};
    % Eligibility
    \node[box]       (elig)  at (0,-7.2)   {Full-text assessment ($n = 496$)};
    \node[eliminated](excl2) at (5.8,-7.2) {Excluded on full text ($n = 337$; EC1--EC4)};
    % Included
    \node[box] (incl) at (0,-9.0) {\textbf{Included in SLR} ($n = \mathbf{159}$)};

    \draw[arr] (db)     -- (dedup);
    \draw[arr] (snow)   -- (dedup);
    \draw[arr] (dedup)  -- (screen);
    \draw[arr] (screen) -- (elig);
    \draw[arr] (screen) -- (excl1);
    \draw[arr] (elig)   -- (incl);
    \draw[arr] (elig)   -- (excl2);
  \end{tikzpicture}
  \caption{PRISMA 2020 flow diagram. The complete checklist is in Online
           Supplement~A.}
  \label{fig:prisma}
\end{figure}

\subsection{Data Extraction}
\label{subsec:extraction}

For each of the 159 included studies, two independent reviewers extracted data
on: (1) the repository representation paradigm(s) employed, (2) the SE
objective(s) targeted, (3) the 10 structural dimensions listed in
Section~\ref{subsec:dimensions}, (4) the LLM architecture used, (5) the
benchmark and evaluation metric, and (6) key quantitative results.
Inter-rater agreement was measured using Cohen's $\kappa$ on paradigm
assignment; $\kappa = 0.82$ (almost perfect agreement~\cite{cruzes2011recommended}).
Disagreements were resolved through discussion and, where necessary, by
consulting a third reviewer.

\subsection{Thematic Synthesis}
\label{subsec:synthesis}

We applied the three-stage thematic synthesis procedure of Cruzes and
Dyb{\aa}~\cite{cruzes2011recommended}: (1) \emph{open coding} of extracted data
into primary codes; (2) \emph{axial coding} to identify relationships between
codes and group them into categories; and (3) \emph{selective coding} to
develop the seven paradigms as higher-order themes.
The thematic synthesis was conducted iteratively over four rounds, with the
taxonomy stabilising after round three (no new paradigms emerged in round four).
